% Copyright (c) 2012-2026 Danny Lee Eldridge. Licensed under the Hippocratic License 3.0 with the Military Activities and Mass Surveillance modules (HL3-MIL-SV), https://firstdonoharm.dev/version/3/0/mil-sv.html
\documentclass[11pt]{article}
\usepackage[utf8]{inputenc}
\usepackage{amsmath,amssymb,amsfonts}
\usepackage[a4paper,margin=1in]{geometry}
\usepackage{booktabs}
\usepackage{hyperref}
\newcommand{\Mp}{M_{\rm Pl}}

\title{\textbf{The Barbero--Immirzi Field as a Hyperbolic Axion--Dilaton:\\
a Fixed Kinetic Mixing $\lambda=-2\sqrt{2/3}$ and the Spiral Twist Hypothesis}}
\author{D.~L.~Eldridge\thanks{Independent researcher, Oklahoma, USA. Contact: dannyeldridge54@gmail.com.
Scripts and data for every number: \url{https://github.com/dannyeldridge54/the-twist-torsion-cosmology};
companion paper and fits: doi:\href{https://doi.org/10.5281/zenodo.23867116}{10.5281/zenodo.23867116}.}}
\date{Theory note, 2 November 2026}

\begin{document}
\maketitle

\begin{abstract}
Promoting the Barbero--Immirzi (BI) parameter to a field in $\mathcal N=1$ supergravity is equivalent to adding a
dilaton--axion chiral superfield with a fixed K\"ahler potential (Gates, Ketov \& Yunes 2009). We point out that
this K\"ahler potential has metric $K_{\Phi\bar\Phi}=3/(\Phi+\bar\Phi)^2$, i.e.\ a hyperbolic field space of
curvature $-2/3$ (the $\alpha=1$ attractor geometry), so the BI axion's kinetic term is multiplied by
$e^{\lambda\chi}$ of the canonical dilaton with $\lambda=-2\sqrt{2/3}\simeq-1.633$ \emph{fixed, with no free
parameter}. This is the kinetic-mixing structure shown to produce an apparent (``illusory'') phantom crossing in
DESI-type analyses while obeying the null energy condition (Toomey et al.\ 2025), whose fit gives
$\lambda=-1.3^{+1.3}_{-1.5}$. We are not aware of earlier work deriving this geometry, and hence $\lambda$, from the
BI field. We then report what survives four consistency checks: (i) with matter on the supergravity frame, the
dilaton has a general-relativity point at $\phi=-\tfrac12$ that is a cosmological attractor, but Cassini forces the
dilaton within $|\Delta\chi|<0.0072\,\Mp$ of it today, which switches the mechanism off; scale (Weyl) invariance
offers a known route to matter blindness; (ii) the simplest supergravity potential (constant superpotential) has
a turning point exactly at $\phi=-\tfrac12$ but negative energy; (iii) the same $\alpha=1$ geometry as an inflaton
gives $n_s=0.967$ at $N=60$, $2.4\sigma$ below ACT DR6; (iv) the dilaton cannot be the early-dark-energy field.
A first background calculation with the decay constant fixed by the Nieh--Yan period ($F=(4/\pi)\sqrt{3/2}\,\Mp$)
behaves as a thawing dilaton ($w_0\simeq-0.84$, $w_a\simeq-0.24$); a literature cross-check of the illusion is not
yet reproduced. We state the model (``Spiral Twist'') and the tests that can falsify it.
\end{abstract}

\section{Status in one paragraph}
This note fixes a theoretical result and a hypothesis, with their checks, before the hypothesis is fitted to data.
The algebra in \S\ref{sec:geometry} is exact and was verified symbolically. Everything after it is a
\emph{hypothesis}: some checks fail or are in tension, and they are reported as such. The companion analysis
(doi:10.5281/zenodo.23867116 and its updates) finds that the minimal late-time Twist field's best preference in
data (DES Y5 supernovae) is a calibration effect: it vanishes when the supernova sample is split at $z=0.1$, and
weakens on DES-Dovekie and Union3. The present note is therefore motivated by theory, not by a fit.

\section{The BI field's supergravity geometry}\label{sec:geometry}
Gates, Ketov \& Yunes \cite{gky2009} show that a dynamical BI parameter in $\mathcal N=1$ supergravity is
equivalent to a chiral superfield $\Phi=\phi+ia$ whose imaginary part couples to the (supersymmetric) Holst term and
whose real part couples to the curvature, with (their eqs.~5.36--5.37)
\begin{equation}
\Omega=e^{\Phi+\bar\Phi}\,(\Phi+\bar\Phi),\qquad K=-3\ln(-\Omega/3),\qquad \phi<0 .
\end{equation}
The linear piece $-3(\Phi+\bar\Phi)$ of $K$ drops out of the metric, leaving
\begin{equation}
K_{\Phi\bar\Phi}=\frac{3}{(\Phi+\bar\Phi)^2},\qquad
\mathcal L_{\rm kin}=-\frac{3}{4\phi^2}\left[(\partial\phi)^2+(\partial a)^2\right]
\end{equation}
(reduced Planck units): the upper-half-plane metric with Gaussian curvature $-2/3$, i.e.\ the $\alpha=1$ case of
$K=-3\alpha\ln(\Phi+\bar\Phi)$. With the canonical dilaton $|\phi|=e^{\sqrt{2/3}\,\chi}$ and canonical axion
$\tilde a=\sqrt{3/2}\,a$,
\begin{equation}
\mathcal L_{\rm kin}=-\tfrac12(\partial\chi)^2-\tfrac12\,e^{\lambda\chi}(\partial\tilde a)^2,\qquad
\boxed{\lambda=-2\sqrt{2/3}\simeq-1.633}.
\end{equation}
The same number follows from the general $\alpha$-attractor relation $\lambda=-2\sqrt{2/(3\alpha)}$
\cite{qalpha2026} at $\alpha=1$. Toomey et al.\ \cite{toomey2025} fit this kinetic coupling to Planck + DESI DR2
and find $\lambda=-1.3^{+1.3}_{-1.5}$ (prior $(-3,0)$); the fixed value lies inside their $1\sigma$ range. This match
was found after their fit and is a consistency, not a test. It becomes a test when $\lambda$ is held fixed in a fit.
The sign of $\lambda$ depends on the orientation of the dilaton.

\section{Matter coupling and the solar system}
If matter couples to the frame in which $\Omega$ multiplies $R$ (``reading A''), the Einstein-frame metric is
$g_E=(-\Omega/3)g_J$ and matter feels $A^2=-3/\Omega$, so
\begin{equation}
\alpha_{\rm DEF}(\chi)=\sqrt2\,\frac{d\ln A}{d\chi}=\frac{2}{\sqrt3}\left(|\phi|-\tfrac12\right).
\end{equation}
The coupling vanishes at $\phi=-\tfrac12$, where $\Omega$ is stationary, and $d^2\ln A/d\chi^2=+1/3$ there: a
Damour--Nordtvedt attractor. Cassini ($|\gamma-1|<2.3\times10^{-5}$, $|\alpha|<3.39\times10^{-3}$) requires
$|\Delta\chi|<0.0072\,\Mp$ today. A dilaton pinned this tightly cannot drive the kinetic-mixing effect, so in
reading A the model reduces to a single axion today. This parallels the Brax--Martin no-go theorem for no-scale quintessence
\cite{brax2007}. The BI field's own kinetic structure is discussed in \cite{taveras2008,calcagni2009}.

\emph{Reading B} (matter blind to the dilaton) is natural if the theory is scale (Weyl) invariant: the dilaton is
then the Goldstone mode of broken scale symmetry with only derivative couplings \cite{blas2011,ferreira2017}, and
Einstein--Cartan gravity admits this symmetry with non-propagating torsion \cite{karananas2021,piani2022}. Whether
reading B can coexist with the $\alpha=1$ geometry above is open.

\section{The simplest supergravity potential}
With a constant superpotential $W$, $V=e^K(K^{\Phi\bar\Phi}|K_\Phi W|^2-3|W|^2)$ gives
\begin{equation}
V(\phi)=\frac{81}{2}\,W^2\,\frac{(-\phi-1)\,e^{-6\phi}}{\phi^2},
\end{equation}
with extrema at $\phi=-\tfrac12$ (a local maximum, exactly at the GR point) and $\phi=-\tfrac23$ (the minimum),
both at negative energy; $V=0$ at $\phi=-1$; for $\phi<-1$ it is positive and steep
($-d\ln V/d\chi$ from 8 to 14), as the illusion requires. It needs an uplift to describe dark energy, and its minimum
is not the GR point.

\section{The Spiral Twist hypothesis}
In the Einstein frame with reading B:
\begin{equation}
\mathcal L=-\tfrac12(\partial\chi)^2-\tfrac12e^{\lambda\chi}(\partial\theta)^2-V_0e^{-\kappa\chi}
-m^2F^2\left[1-\cos(\theta/F)\right],
\end{equation}
with $\lambda=-2\sqrt{2/3}$ (above) and $F=(4/\pi)\sqrt{3/2}\,\Mp\simeq1.56\,\Mp$, the value at which the
Nieh--Yan index has period one for a Planck-length torsion scale (companion analysis). The angle $\theta$ is the BI
(Twist) field and the radius $\chi$ grows as the spiral unwinds. Free: $\kappa$, $m$, $\theta_i$; $V_0$ is set by
closure.

\paragraph{First background pass.} A two-field solver (validated: exact $\Lambda$CDM limit; single-field limit;
null energy condition) with a CPL fit to $H(z)$ gives, over 80 settings, at best $w_0\simeq-0.84$, $w_a\simeq-0.24$
(at $\kappa=1$), the thawing relation $w_a\simeq-1.53(1+w_0)$; with $F=1.56\,\Mp$ the angle carries too little
energy for the mixing to matter. Using Toomey et al.'s effective dilaton equation of state, the illusion appears and
is stronger with the fixed $\lambda$ ($w_{\rm eff}(0)=-1.50$ against $-1.09$ at their $\lambda$), but in the opposite
time direction to their posterior; their benchmark inputs are not fully specified, and this cross-check is
\textbf{not yet reproduced}. These numbers are provisional.

\section{Consistency checks (rules fixed before computing)}
\begin{center}\small
\begin{tabular}{p{0.34\textwidth}p{0.42\textwidth}p{0.16\textwidth}}
\toprule
Check & Result & Verdict \\
\midrule
Inflation from $\alpha=1$ ($n_s=1-2/N$, $r=12/N^2$) & $n_s=0.9667$ at $N=60$ vs $0.974\pm0.003$ & tension ($2.4\sigma$) \\
Dilaton relaxation as early dark energy & $\Omega_\chi(z{=}3500)\le0.025$; $G$ changes $\times13$ at BBN & too small \\
Solar system (reading A) & attractor at $\phi=-\tfrac12$, band $0.0072\,\Mp$ & passes only if pinned \\
\bottomrule
\end{tabular}
\end{center}

\section{How to falsify it}
\begin{enumerate}
\item Fit the Spiral Twist with $\lambda$ and $F$ fixed (two dark-energy parameters, as in $w_0$--$w_a$) to CMB, DESI
and the three supernova samples; it fails if it does no better than the thawing single field.
\item Derive reading B from a Weyl-invariant Einstein--Cartan action containing the BI field; if the dilaton
necessarily couples to matter, the model is excluded by Cassini.
\item With reading B, the dilaton is a light field coupled only to the dark sector: parity-violating gravitational
waves and CMB birefringence from the angle remain the direct probes (companion analysis).
\end{enumerate}

\section*{Acknowledgements}
The author used AI assistants (Claude, Anthropic) for algebra checks, code and drafting; every number is reproduced
by scripts in the repository (\texttt{bi\_kahler\_lambda.py}, \texttt{solar\_lowz.py}, \texttt{bi\_superpotential.py},
\texttt{ufe\_checks.py}, \texttt{spiral\_w.py}), and the author takes full responsibility for the content.

\begin{thebibliography}{99}
\bibitem{gky2009} S.~J.~Gates, S.~V.~Ketov, N.~Yunes, Phys.\ Rev.\ D 80, 065003 (2009), arXiv:0906.4978.
\bibitem{toomey2025} M.~W.~Toomey, E.~Hughes, M.~M.~Ivanov, J.~M.~Sullivan, arXiv:2511.23463.
\bibitem{qalpha2026} \emph{Quintessential $\alpha$-attractors, updated}, arXiv:2608.19781.
\bibitem{brax2007} P.~Brax, J.~Martin, Phys.\ Lett.\ B 647, 320 (2007), hep-th/0612208.
\bibitem{blas2011} D.~Blas, M.~Shaposhnikov, D.~Zenh\"ausern, arXiv:1104.1392.
\bibitem{ferreira2017} P.~G.~Ferreira, C.~T.~Hill, G.~G.~Ross, ``No fifth force in a scale invariant universe'',
Phys.\ Rev.\ D 95, 064038 (2017).
\bibitem{karananas2021} G.~K.~Karananas, M.~Shaposhnikov, A.~Shkerin, S.~Zell, arXiv:2108.05897.
\bibitem{piani2022} M.~Piani, J.~Rubio, JCAP 05 (2022) 009, arXiv:2202.04665.
\bibitem{taveras2008} V.~Taveras, N.~Yunes, arXiv:0807.2652.
\bibitem{calcagni2009} G.~Calcagni, S.~Mercuri, arXiv:0902.0957.
\bibitem{act2025} T.~Louis et al.\ (ACT), arXiv:2503.14452.
\end{thebibliography}
\end{document}
